\documentclass[11pt,a4paper]{article}
\usepackage[margin=2.2cm]{geometry}
\usepackage{amsmath,amssymb}
\usepackage{graphicx}
\usepackage{booktabs}
\usepackage{xcolor}
\usepackage{hyperref}
\usepackage{caption}
\usepackage[section]{placeins}
\graphicspath{{figs/}}

\title{Local optimisation of the SHiP muon shield with per-muon surrogates:\\[0.3ex]
\large the LCSO surrogates in a trust-region loop on the MuonShieldChallenge}
\author{}
\date{\today}

\begin{document}
\maketitle

\begin{abstract}
The per-muon hit surrogates of this repository (DeepONet, full and reduced Taylor, with and
without residual) are used as the model of a trust-region optimiser that retrains the surrogate
around the current design after every step, and are compared on the two problems of the
MuonShieldChallenge at a fixed budget of simulated muons. Four choices make the difference: designs
sampled in a box wider than the step (radius $0.05$ of the bounds, against steps of
$0.005$--$0.04$), the per-muon hits of every full simulation reused as training data, an acceptance
test that allows for the Poisson noise of the hit count, and the full Taylor surrogate (no
residual) as the model. Letting the trust radius grow beyond the sampling box (the surrogate only supplies
the gradient and curvature) helps on hard by a further 25\%. On the hard problem (72 parameters, $5\times10^8$ muons per simulation,
noisy) the result reaches 647 hits (median of 2 seeds, from $\sim$69\,500) with 100 full
simulations' worth of muons, against 3220 for the best baseline (TuRBO) at the same budget and
1380 for TuRBO at five times the budget. On the easy problem (30 parameters, $10^7$ muons) it
reaches 36 hits at budget 50 and 18 at 400 (median of 5 seeds, from 1173), ahead of every
black-box baseline (45--75 at 400) and of the earlier global LCSO (58), but behind the RL
methods of the challenge (5--13 at 400).
\end{abstract}

\tableofcontents

%==============================================================================
\section{The problems and the budget}
%==============================================================================

The MuonShieldChallenge defines two problems, both minimising the number of muons reaching the
sensitive plane over the parameters $u\in[0,1]^D$ (scaled to the bounds) under exact length,
cost and (hard) cavern constraints:
\begin{itemize}
  \item \textbf{easy}: $D=30$, $10^7$ muons per full simulation, fixed simulation seed
        (deterministic), reference design 1173 hits;
  \item \textbf{hard}: $D=72$, $5\times10^8$ muons per full simulation, a new seed every
        simulation (noisy, $\sigma\approx\sqrt f$), cavern walls, reference $\approx69\,500$ hits.
\end{itemize}
The budget counts simulated muons, in units of full simulations: simulating 32 designs on
$2.5\times10^5$ muons each costs $0.8$ on easy. Only full simulations give a best $f$, which is the
figure of merit at every budget. Infeasible designs are never simulated.

Two properties of the simulator set what an optimiser can do:
\begin{itemize}
  \item \emph{A geometry change reshuffles each muon's scattering}: the per-muon hits of two nearby
        designs are nearly independent draws, even with a fixed seed. The difference of $f$
        between two designs therefore carries a noise of $\sqrt{f_1+f_2}$, also on the
        deterministic easy problem, and common random numbers (the same muons for all designs)
        remove little of it.
  \item \emph{Hits cannot be predicted from the muons that hit at the reference}: 10--25\% of the
        hits at nearby designs come from muons that did not even reach the plane at the reference,
        so simulating only ``near-miss'' muons is not an option.
\end{itemize}

%==============================================================================
\section{The optimiser}
%==============================================================================

\texttt{baselines/lcso\_local.py} in the challenge (method \texttt{lcso\_local}) imports the
surrogate architectures of \texttt{lcso/models.py} from this repository. Each round:
\begin{enumerate}
  \item \textbf{Sample} $n_d=32$ feasible designs in the box $|u-u_\text{inc}|_\infty\le r_s$ around the
        incumbent (the incumbent first) and simulate each on $m$ muons: $2.5\%$ of the sample
        ($2.5\times10^5$ on easy, $1.25\times10^7$ on hard). Keep every hit and a fraction $k$ of the
        misses (case-control; the training logit carries the offset $-\log k$).
  \item \textbf{Train} the surrogate $p(\text{hit}\mid u,x)$ on the rows of the last $B$ rounds, in
        coordinates centred on the incumbent, $v=(u-u_\text{inc})/r_s$. A Taylor surrogate therefore
        expands around the current point, and its gradient and Hessian there are explicit
        coefficients. The model is warm-started from the previous round (500 Adam steps; 3000 in
        the first round).
  \item \textbf{Step} on the predicted hits $F(u)=\tfrac{N}{n}\sum_{x}\sigma(\text{logit}(u,x))$
        over a fixed random set of $n$ muons ($2\times10^6$ on easy, $5\times10^6$ on hard):
        a trust-region step of radius $r_t$ within the bounds and the exact constraints, either
        on the linear model (\texttt{gd}), on the quadratic model from the autograd gradient and
        Hessian (\texttt{newton}), or minimising $F$ itself with SLSQP (\texttt{surrogate}).
  \item \textbf{Accept} after a full simulation of the candidate. With the noise-aware test, the step
        is rejected only if the candidate is worse by more than $\tau\sigma$,
        $\sigma=\sqrt{f_\text{inc}+f_\text{cand}}$; the trust radius halves only then and doubles if the
        actual reduction exceeds a fraction $g$ of the predicted one at the boundary
        ($0.005\le r_t\le r_\text{max}$; $g=0.75$, $r_\text{max}=0.2$ unless stated).
  \item \textbf{Reuse}: the candidate's full simulation is also training data. Its hits and a
        fraction $k$ of its misses join the next round's rows.
\end{enumerate}
The final configuration, \texttt{C\_follow2}, is the full Taylor surrogate (\texttt{kind='taylor'}:
the full gradient and a full-rank curvature, no residual), \texttt{surrogate} steps, $r_s=0.05$
growing as $\max(0.05,\,2r_t)$, $r_t=0.02$ initially, $k=0.05$, $B=5$, $\tau=1$, with reuse. The
hard runs used an earlier configuration (Newton steps, $k=0.02$, $B=3$, fixed $r_s=0.05$): the easy
sweeps that improved on it finished after the hard runs had started. The best hard configuration,
\texttt{HF\_g025}, adds to it a trust radius free to leave the sampling box ($g=0.25$,
$r_\text{max}=1$, Section~\ref{sec:far}).

%==============================================================================
\section{What mattered: the easy sweeps}
%==============================================================================

All settings were tuned on easy, at budget 200 (3 seeds) and then 400 (3--5 seeds), one change
at a time from the current best. Table~\ref{tab:steps} follows the path from the first version
(210 hits at budget 200) to the final one (21); Appendix~\ref{app:sweeps} lists every setting
tried.

\begin{table}[h]\centering\small
\begin{tabular}{lrrrrr}
\toprule
best hits (median over seeds) at budget & 10 & 25 & 50 & 100 & 200 \\
\midrule
first version: reduced Taylor $r{=}8$ + residual, $r_s=0.02$, Newton (3) & 515 & 418 & 334 & 257 & 210 \\
$+$ sampling radius $r_s=0.05$ (3) & 378 & 242 & 216 & 90 & 65 \\
$+$ noise-aware acceptance, $\tau=1$ (3) & 378 & 267 & 136 & 69 & 60 \\
$+$ full simulations reused as data (3) & 444 & 95 & 63 & 56 & 49 \\
$+$ full Taylor surrogate (5) & 221 & 72 & 41 & 31 & 31 \\
$+$ \texttt{surrogate} steps, $k=0.05$, $B=5$ (5) & 103 & 61 & 46 & 33 & 23 \\
$+$ sampling radius following the trust radius (5, final) & 136 & 61 & 36 & 24 & 21 \\
\bottomrule
\end{tabular}
\caption{Easy: the changes that helped, each added to the previous line (number of seeds in
brackets).}
\label{tab:steps}
\end{table}

\paragraph{Sampling radius.} The single largest gain. With $r_s=0.02$ the designs of a round differ
by a few hits, comparable with the noise of a design simulated on $2.5\times10^5$ muons, and
the trust radius collapses to its floor. Too small is much worse than too large: $r_s=0.01$ ends at
299, $0.02$ at 210, $0.05$ at 65, $0.08$ at 114, $0.1$ at 121 (budget 200). Tying the sampling box
to the trust radius (\texttt{couple}) brings back the collapse (308). Letting it \emph{grow} with the
trust radius but never shrink below $0.05$ helps a little (the last line of
Table~\ref{tab:steps}). Note that this is much smaller than the $\pm10\%$ box of the robustness
analysis: for optimisation only the local model matters.

\paragraph{The noise floor.} Even with a good model, a step of $r_t=0.005$ predicts reductions of
5--20 hits at $f\approx200$, below the noise $\sqrt{2f}\approx20$. A strict ratio test
($\rho=\text{actual}/\text{predicted}\ge0.1$) then rejects most good steps, shrinks the radius, and
the run stalls (the first version spent 60 rounds at $f=257$). The noise-aware test accepts a
step unless it is significantly worse, so the incumbent keeps moving along the predicted descent.
$\tau=1$ and $\tau=2$ give the same result.

\paragraph{Reusing the full simulations.} Every round pays a full simulation for the acceptance
test: $10^7$ muons at the candidate, more than the $8\times10^6$ of the 32 sampled designs.
Keeping its per-muon hits as training data costs nothing and halved the hits at budget 25.

\paragraph{The muons per design and the number of designs.} With the budget fixed, $2.5\times10^5$
muons per design (2.5\% of the sample) beats $5\times10^4$, $10^5$ and $10^6$; 32 designs per
round beat 16 and 64. Fewer designs per round means more steps, but the model needs about $D$
designs to see every direction. Sampling new designs only every 2--4 rounds (training in between
on the reused full simulations alone) did not help.

\paragraph{Training.} Warm-starting from the previous round, cold restarts, a buffer of 1--8 rounds
and 200--1500 steps per round all end within the scatter of the seeds. Keeping 5\% of the misses
instead of 2\% helps slightly.

\paragraph{What did not help.}
\begin{itemize}
  \item \emph{Common random numbers} (the same muons for every design of a round): no consistent
        effect. Without them the first version is better (110 against 210 at budget 200), the
        $r_s=0.05$ one worse (100 against 65). The geometry reshuffles the scattering, so a shared
        muon set removes little of the noise of the comparisons.
  \item \emph{Importance sampling of the muons} by the surrogate's hit probability (with weights
        $1/(Nq)$ in the loss), and \emph{acceptance on a subset} of $10^6$ muons with a full
        simulation every third accepted step. Both give cheaper rounds but noisier models and
        decisions: at best equal (\texttt{is05\_eval1e6}: 103 against 65 for the same base).
  \item \emph{Gradient steps}: gd ends at 82 against 65 for Newton (budget 200), 27 against 24 with
        the full Taylor surrogate (budget 400). Minimising the whole surrogate (\texttt{surrogate})
        is the best start (103--127 at budget 10) and ends level with Newton.
\end{itemize}

%==============================================================================
\section{Which surrogate}
%==============================================================================

\begin{figure}[h]\centering
\includegraphics[width=0.8\linewidth]{easy_kinds.png}
\caption{Easy: the surrogate kinds in the same optimiser (sampling radius 0.05, noise-aware
acceptance, reuse, Newton steps). Median over 5 seeds and the range of the seeds.}
\label{fig:kinds}
\end{figure}

\begin{table}[h]\centering\footnotesize
\begin{tabular}{lrrrrrr}
\toprule
best hits at budget & 10 & 25 & 50 & 100 & 200 & 400 \\
\midrule
\multicolumn{7}{l}{\emph{easy, median of 5 seeds}} \\
full Taylor & 221 & 72 & 41 & 31 & 31 & 24 \\
reduced Taylor $r=16$ + residual & 200 & 67 & 58 & 55 & 43 & 29 \\
DeepONet & 138 & 65 & 51 & 41 & 34 & 34 \\
reduced Taylor $r=8$ + residual & 444 & 125 & 105 & 59 & 45 & 34 \\
reduced Taylor $r=8$ & 240 & 78 & 61 & 57 & 37 & 35 \\
\multicolumn{7}{l}{\emph{easy, budget 200 only, 3 seeds, without reuse}} \\
reduced Taylor $r=2$ / $r=4$ + residual & 576 / 467 & 576 / 367 & 405 / 132 & 379 / 114 & 355 / 114 & -- \\
full Taylor + residual & 230 & 127 & 101 & 96 & 96 & -- \\
\midrule
\multicolumn{7}{l}{\emph{hard, median of 2 seeds}} \\
full Taylor & 14\,520 & 3892 & 1034 & 859 & -- & -- \\
DeepONet & 18\,540 & 7443 & 1976 & 1226 & -- & -- \\
reduced Taylor $r=8$ + residual & 39\,570 & 30\,980 & 25\,770 & 14\,130 & -- & -- \\
\bottomrule
\end{tabular}
\caption{The surrogate kinds in the same optimiser (hard: column 100 is the end of the runs,
budget 99).}
\label{tab:kinds}
\end{table}

The full Taylor surrogate (a full gradient and a full-rank curvature around the incumbent, no
residual) is the best model for this loop on both problems (Table~\ref{tab:kinds},
Fig.~\ref{fig:kinds}). The reason is the shape of the loop: the surrogate only has to be right
within $r_s=0.05$ of the incumbent, where a quadratic logit is enough, and the step only uses its
value, gradient and curvature there. The residual adds capacity that this small box does not need
and data that 32 designs per round cannot pin down: with the residual the full Taylor surrogate
ends at 96 instead of 77 (budget 200, without reuse).

The rank of the reduced Taylor surrogate matters a lot. Rank 2 and 4 are too few for 30
parameters (355 and 114). Rank 8 works on easy but not on hard, where with 72 parameters it is
the only kind that fails (20\,192 and 7897 hits, one seed stuck for 35 rounds on a plateau at
40\,000--48\,000). Rank 16 is as good as the full one at budget 400 on easy. The DeepONet is the best
start (138 at budget 10), as on the toy problems of the comparison report, where it reached a
moderate accuracy with the fewest samples. It then falls behind the full Taylor surrogate, and is
second on hard.

%==============================================================================
\section{Results against the baselines}
%==============================================================================

\begin{figure}[h]\centering
\includegraphics[width=0.85\linewidth]{hard_curves.png}
\caption{Hard: best hits against the muons simulated. Median over the seeds (number in brackets);
 the DeepONet and reduced Taylor runs use $g=0.75$, $r_\text{max}=0.2$;
the bands are the range of the local LCSO seeds.}
\label{fig:hard}
\end{figure}

\begin{figure}[h]\centering
\includegraphics[width=0.85\linewidth]{easy_curves.png}
\caption{Easy: best hits against the muons simulated, for the final configuration, the same
optimiser with the DeepONet, the first version, and the challenge's baselines.}
\label{fig:easy}
\end{figure}

\begin{table}[h]\centering\small
\begin{tabular}{lrrrrrr}
\toprule
best hits (median over seeds) at budget & 10 & 25 & 50 & 75 & 100 & 500 \\
\midrule
\textbf{local LCSO, full Taylor, far steps} (2) & \textbf{13\,060} & \textbf{4010} & \textbf{1010} & \textbf{675} & \textbf{647} & -- \\
local LCSO, DeepONet (2) & 18\,540 & 7443 & 1976 & 1352 & 1226 & -- \\
TuRBO (3) & 55\,350 & 21\,450 & 9710 & 5391 & 3220 & 1380 \\
LCSO global (old), Newton (3) & 69\,540 & 37\,490 & 21\,590 & 9177 & 5093 & 2933 \\
BO (3) & 69\,640 & 33\,750 & 33\,750 & 12\,520 & 8160 & 2669 \\
CMA-ES (3) & 69\,210 & 69\,210 & 69\,200 & 36\,770 & 36\,770 & 3201 \\
GA (3) & 69\,620 & 69\,500 & 46\,650 & 46\,650 & 26\,470 & 3517 \\
\bottomrule
\end{tabular}
\caption{Hard. The local LCSO runs end at budget 100 (shown at 99).}
\label{tab:hard}
\end{table}

\paragraph{Hard.} Local LCSO with the full Taylor surrogate and far steps is the best method at
every budget (Table~\ref{tab:hard}, Fig.~\ref{fig:hard}). At budget 25 it is where TuRBO is at
100. At budget 50 it already beats every baseline at budget 500. Its two seeds end at 619 and 675
hits. With the DeepONet in the same loop (and $g=0.75$, $r_\text{max}=0.2$) the runs end at 944 and
1509. The hard runs used Newton steps; the easy setting with \texttt{surrogate} steps and a
growing sampling box was not run on hard. Each hard run takes 3.5--5 hours on one GPU.

\begin{table}[h]\centering\small
\begin{tabular}{lrrrrrr}
\toprule
best hits (median over seeds) at budget & 10 & 25 & 50 & 100 & 200 & 400 \\
\midrule
RL flow (3) & 60 & 33 & 19 & 10 & 5 & 5 \\
RL, model-based (3) & 176 & 43 & 22 & 22 & 20 & 13 \\
\textbf{local LCSO, full Taylor, final} (5) & \textbf{136} & \textbf{61} & \textbf{36} & \textbf{24} & \textbf{21} & \textbf{18} \\
local LCSO, DeepONet (5) & 138 & 65 & 51 & 41 & 34 & 34 \\
GA (3) & 1173 & 411 & 156 & 126 & 85 & 45 \\
CMA-ES (3) & 1163 & 390 & 390 & 164 & 83 & 53 \\
LCSO global (old), Newton (2) & 892 & 112 & 96 & 96 & 92 & 58 \\
TuRBO (3) & 1173 & 1085 & 317 & 130 & 72 & 67 \\
BO (3) & 1173 & 750 & 425 & 141 & 75 & 75 \\
\bottomrule
\end{tabular}
\caption{Easy. The black-box baselines whose runs end just short of budget 400 are shown at 399.}
\label{tab:easy}
\end{table}

\paragraph{Easy.} The final configuration beats every black-box baseline and the earlier global
LCSO at every budget (Table~\ref{tab:easy}, Fig.~\ref{fig:easy}): 36 hits at budget 50, where
TuRBO, CMA-ES, GA and BO are at 156--425, and 18 at 400 against 45--75. The RL methods of the
challenge are better: RL flow reaches 19 at 50 and 5 at 200, the model-based RL 22 at 50 and 13 at
400. They learn a policy over designs with partial simulations, and are not local.
The relative gain of local LCSO is larger on hard (more parameters, more muons per simulation),
where the RL methods have not been run.

%==============================================================================
\section{Letting the step leave the training box}\label{sec:far}
%==============================================================================

The surrogate is trained on designs within $r_s=0.05$ of the incumbent, but a Newton step only uses
its gradient and Hessian there. Whether the step is good is decided by the full simulation, so
nothing prevents the trust radius from growing beyond the training box when the steps keep paying
off. Two parameters control this: the cap $r_\text{max}$ and the growth threshold $g$ (the radius
doubles when a step at the boundary achieves at least $g$ of the predicted reduction).

\begin{table}[h]\centering\small
\begin{tabular}{lrrrrrr}
\toprule
best hits (median over seeds) at budget & 10 & 25 & 50 & 75/100 & 99/200 & 400 \\
\midrule
\multicolumn{7}{l}{\emph{easy, Newton steps}} \\
$g=0.75$, $r_\text{max}=0.2$ (10) & 138 & 65 & 47 & 34 & 25 & 21 \\
$g=0.75$, $r_\text{max}=1$ (5) & \multicolumn{6}{c}{identical to the line above: $r_t$ never reaches 0.2} \\
$g=0.5$, $r_\text{max}=1$ (5) & 183 & 81 & 48 & 29 & 22 & 17 \\
$g=0.5$, $r_\text{max}=1$, $r_t=0.1$ initially (5) & 182 & 62 & 45 & 30 & 26 & 21 \\
$g=0.25$, $r_\text{max}=1$ (10) & 170 & 59 & 40 & 32 & 22 & 19 \\
\multicolumn{7}{l}{\emph{easy, final configuration (\texttt{surrogate} steps, following box)}} \\
$g=0.75$, $r_\text{max}=0.2$ (5) & 136 & 61 & 36 & 24 & 21 & 18 \\
$g=0.25$, $r_\text{max}=1$ (5) & 138 & 77 & 55 & 36 & 29 & 21 \\
\midrule
\multicolumn{7}{l}{\emph{hard, Newton steps (columns 75 and 99)}} \\
$g=0.75$, $r_\text{max}=0.2$ (2) & 14\,520 & 3892 & 1034 & 889 & 859 & -- \\
$g=0.5$, $r_\text{max}=1$ (2) & 17\,610 & 4572 & 1269 & 824 & 735 & -- \\
$g=0.25$, $r_\text{max}=1$ (2) & 13\,060 & 4010 & 1010 & 675 & 647 & -- \\
\bottomrule
\end{tabular}
\caption{A trust radius free to grow beyond the sampling box $r_s=0.05$. Easy columns are budgets
10--400; hard columns 10, 25, 50, 75, 99.}
\label{tab:far}
\end{table}

Table~\ref{tab:far} gives the results.
\begin{itemize}
  \item \emph{The cap was never the limit.} With the default $g=0.75$ the radius never reaches 0.2,
        and raising the cap changes nothing. What keeps the steps short is the growth test.
  \item \emph{With a lower threshold the steps do leave the box,} but not far. With $g=0.25$
        every easy seed reaches $r_t=0.16$. Between 4 and 23 steps per run go beyond $0.05$, and
        about half of them improve $f$. On hard with $g=0.5$, one seed takes 18 of its 52 steps at
        $0.08$, and 15 of them improve. No run goes beyond 0.2: away from the box the model
        overpredicts the reduction (one hard step predicted $-844$ hits and gave $+49$), so the
        radius stops doubling.
  \item \emph{Hard: a consistent gain.} All four hard runs with $r_\text{max}=1$ (619--753) end below
        both runs without it (767, 950), a 25\% lower median with $g=0.25$.
  \item \emph{Easy: within the scatter.} With 10 seeds, Newton with $g=0.25$ ends at 19 against 21
        (means 17.7 and 20.3), but starts worse (170 against 138 at budget 10). On top of the final
        configuration it does not help (21 against 18): there the sampling box already follows the
        trust radius, $r_s=\max(0.05,2r_t)$, so a longer step also widens the training box and
        thins out the data near the incumbent.
\end{itemize}
On easy, the runs end in the noise-dominated regime of Section~\ref{sec:disc}, where longer steps
cannot help much. On hard, the curvature of the surrogate stays useful well beyond the box it was
trained in.

%==============================================================================
\section{Discussion}\label{sec:disc}
%==============================================================================

\paragraph{Why it works.} A per-muon surrogate turns each simulated design into thousands of
labelled rows. 32 designs on $2.5\times10^5$ muons pin down a 30-parameter gradient and
curvature that 32 values of $f$ never could. Retraining around the incumbent keeps the model in
the regime where it is accurate, so the trust region is set by the noise of the simulator, not by
model error. On hard, a design simulated on $1.25\times10^7$ muons has $\sim1700$ hits, and the
model sees the local landscape with little noise.

\paragraph{Where it stops.} On easy the runs end in a noise-dominated regime: at $f\approx20$ the
noise of a comparison is $\sqrt{40}\approx6$ hits, and a step of the minimum radius predicts
reductions of a few hits. A single local run then cannot tell its moves apart. Lower values would
need either more muons per decision (larger budgets) or larger moves (a global component). The RL
methods' advantage on easy probably comes from the latter.

\paragraph{Caveats.} The settings were tuned on easy with 3 seeds and the top ones confirmed with 5;
differences of $\lesssim30\%$ between settings are within the scatter of the seeds (the range of
the seeds is shown in the figures). The hard comparison has 2 seeds per surrogate. On the
deterministic easy problem a lucky design counts as much as a genuinely good one. Every method
benefits from this equally, through the same best-$f$ metric.

%==============================================================================
\section{Reproducing}
%==============================================================================

{\sloppy In the MuonShieldChallenge, with the simulator environment of \texttt{Robustness\_MS/env.sh}
(set \texttt{ROBUSTNESS\_MS\_DIR} if this repository is not next to it):\par}
{\footnotesize
\begin{verbatim}
python baselines/run.py lcso_local configs/easy.json --budget 400 --seed 0 --options \
  '{"kind":"taylor","step":"surrogate","sample_radius":0.05,"follow":2,"keep_neg":0.05,
    "buffer_rounds":5,"noise_tol":1,"reuse_full":true}'
python baselines/run.py lcso_local configs/hard.json --budget 100 --seed 0 --options \
  '{"kind":"taylor","sample_radius":0.05,"noise_tol":1,"reuse_full":true,
    "muons_per_design":12500000,"n_eval":5000000,"trust_max":1,"grow_ratio":0.25}'
# a sweep of option sets x seeds on GPUs 0 and 1, and the table of best f against the budget
python baselines/sweep_lcso_local.py sets.json configs/easy.json --budget 400 --gpus 0 1
python baselines/budget_table.py results/baselines/easy.json results/lcso_local*/*/easy.json
\end{verbatim}}
\texttt{report/optimisation/plots.py} makes the figures. On hard a full simulation is run in
chunks of $5\times10^7$ muons (\texttt{Evaluator.\_hits}): the same distribution of $f$, with
28\,GB of host memory instead of 57.

\appendix
%==============================================================================
\section{Every setting tried on easy}
\label{app:sweeps}
%==============================================================================

Median best hits over 3 seeds unless stated. Changes are relative to the base of each block.

\begin{table}[h]\centering\footnotesize
\resizebox{\linewidth}{!}{\begin{tabular}{llrrrrr}
\toprule
setting & change & @10 & @25 & @50 & @100 & @199 \\
\midrule
\multicolumn{7}{l}{\emph{block 1: base = reduced Taylor $r{=}8$ + residual, Newton, $r_s=r_t=0.02$, 32 designs, $2.5\times10^5$ muons, CRN, $k=0.02$, $B=3$}} \\
base & & 515 & 418 & 334 & 257 & 210 \\
s01 / s05 / s10 & $r_s=0.01$ / $0.05$ / $0.1$ & 794 / 378 / 549 & 566 / 242 / 267 & 524 / 216 / 174 & 299 / 90 / 161 & 299 / 65 / 121 \\
t05 & $r_t=0.05$ & 558 & 416 & 372 & 201 & 140 \\
couple & $r_s=r_t$ & 892 & 593 & 488 & 352 & 308 \\
mpd5e4 / 1e5 / 1e6 & muons per design & 965 / 675 / 1070 & 621 / 463 / 628 & 573 / 330 / 136 & 303 / 291 / 123 & 192 / 227 / 102 \\
nd16 / nd64 & designs per round & 480 / 673 & 283 / 384 & 271 / 317 & 265 / 284 & 262 / 130 \\
first64\_nd16 & 64 then 16 & 523 & 341 & 321 & 293 & 293 \\
nocrn & no CRN & 654 & 233 & 188 & 117 & 110 \\
steps200 / 1500 & training steps & 584 / 575 & 411 / 343 & 229 / 327 & 145 / 307 & 102 / 204 \\
buf1 / buf8 & buffer rounds & 752 / 546 & 534 / 388 & 426 / 258 & 185 / 215 & 154 / 165 \\
cold & no warm start & 520 & 221 & 158 & 152 & 151 \\
gd / surr & step & 516 / 474 & 305 / 312 & 232 / 289 & 213 / 212 & 194 / 149 \\
is05 / is09 & importance 0.5 / 0.9 & 571 / 455 & 338 / 297 & 255 / 274 & 151 / 206 & 108 / 168 \\
eval1e6 & accept on $10^6$ muons & 808 & 340 & 308 & 181 & 108 \\
is05\_eval1e6 & both & 364 & 305 & 213 & 115 & 103 \\
\midrule
\multicolumn{7}{l}{\emph{block 2: base = s05}} \\
nt1 / nt2 & noise-aware, $\tau=1$ / 2 & 378 / 378 & 267 / 218 & 136 / 95 & 69 / 78 & 60 / 53 \\
reuse & full simulations as data & 444 & 95 & 82 & 72 & 72 \\
reuse\_nt1 & both & 444 & 95 & 63 & 56 & 49 \\
s08 & $r_s=0.08$ & 604 & 328 & 162 & 134 & 114 \\
t05 & $r_t=0.05$, floor 0.01 & 428 & 221 & 116 & 75 & 69 \\
nocrn / cold / steps200 & & 475 / 389 / 481 & 214 / 318 / 326 & 165 / 224 / 209 & 103 / 155 / 145 & 100 / 107 / 108 \\
mpd1e6 & & 524 & 315 & 217 & 80 & 75 \\
is05\_eval & & 547 & 186 & 165 & 101 & 89 \\
gd & & 412 & 338 & 159 & 102 & 82 \\
kinds & DeepONet / full Taylor / full Taylor + res. & 196 / 243 / 230 & 88 / 94 / 127 & 74 / 86 / 101 & 74 / 77 / 96 & 67 / 77 / 96 \\
kinds & reduced $r{=}8$ / $r{=}2,4,16$ + res. & 233 / 576, 467, 347 & 144 / 576, 367, 165 & 84 / 405, 132, 144 & 59 / 379, 114, 70 & 59 / 355, 114, 63 \\
\bottomrule
\end{tabular}}
\end{table}

\begin{table}[h]\centering\footnotesize
\resizebox{\linewidth}{!}{\begin{tabular}{llrrrrrr}
\toprule
setting & change & @10 & @25 & @50 & @100 & @200 & @399 \\
\midrule
\multicolumn{8}{l}{\emph{block 3: base R = reuse\_nt1, budget 400}} \\
R (5 seeds) & & 444 & 125 & 105 & 59 & 45 & 34 \\
kinds (5 seeds) & full Taylor & 221 & 72 & 41 & 31 & 31 & 24 \\
 & DeepONet & 138 & 65 & 51 & 41 & 34 & 34 \\
 & reduced $r{=}16$ + res. & 200 & 67 & 58 & 55 & 43 & 29 \\
 & reduced $r{=}8$ & 240 & 78 & 61 & 57 & 37 & 35 \\
nt2 & $\tau=2$ & 444 & 95 & 50 & 50 & 50 & 41 \\
s03 / s07 & $r_s$ & 332 / 360 & 86 / 170 & 79 / 80 & 67 / 50 & 54 / 42 & 54 / 40 \\
t05 & $r_t=0.05$, floor 0.01 & 401 & 182 & 115 & 72 & 59 & 47 \\
nd16 & & 366 & 84 & 77 & 77 & 64 & 52 \\
tmin0025 & floor 0.0025 & 444 & 95 & 63 & 55 & 46 & 40 \\
\midrule
\multicolumn{8}{l}{\emph{block 4: base T = R with the full Taylor surrogate}} \\
kn05 / kn01 & $k=0.05$ / 0.01 & 154 / 250 & 65 / 69 & 52 / 44 & 35 / 26 & 31 / 24 & 20 / 23 \\
buf5 & $B=5$ & 195 & 66 & 45 & 33 & 27 & 22 \\
follow2 & $r_s=\max(0.05,2r_t)$ & 220 & 69 & 42 & 38 & 22 & 22 \\
surr / gd & step & 127 / 218 & 65 / 76 & 37 / 47 & 33 / 37 & 25 / 37 & 22 / 27 \\
nd64 & 64 designs, $1.25\times10^5$ muons & 213 & 89 & 49 & 40 & 32 & 23 \\
t04 & $r_t=0.04$ & 196 & 68 & 51 & 47 & 31 & 31 \\
steps200 & & 169 & 60 & 51 & 46 & 36 & 28 \\
de2 / de3 / de4 & designs every 2 / 3 / 4 rounds & 239 / 306 / 255 & 81 / 90 / 70 & 53 / 53 / 53 & 39 / 34 / 46 & 30 / 29 / 33 & 28 / 21 / 29 \\
de2\_t04 & & 215 & 71 & 42 & 36 & 30 & 22 \\
D\_de2 & DeepONet, de2 & 138 & 78 & 47 & 38 & 32 & 29 \\
\midrule
\multicolumn{8}{l}{\emph{block 5: combinations, 5 seeds}} \\
C & surr, $k=0.05$, $B=5$ & 103 & 61 & 46 & 33 & 23 & 22 \\
C\_newton & Newton, $k=0.05$, $B=5$ & 152 & 63 & 45 & 32 & 25 & 20 \\
C\_follow2 (final) & C + follow2 & 136 & 61 & 36 & 24 & 21 & 18 \\
\bottomrule
\end{tabular}}
\end{table}

\end{document}
